\subsection*{Vídeos do YouTube}
Palestras e tutoriais, a maioria em inglês; os marcados com \pt são em
português.

\small
\begin{longtable}{@{}p{1.8cm} p{7.2cm} p{3.4cm}@{}}
\toprule
Tema & Vídeo & Quem \\
\midrule \endhead
Cobertura & \citetitle{video:aniche-coverage} &
  \citeauthor{video:aniche-coverage} \\
Mutação & \citetitle{video:coles-mutants} (GOTO 2019) &
  \citeauthor{video:coles-mutants}, criador do PIT \\
Mutação & \citetitle{video:pit-maven-tutorial} (tutorial) &
  \citeauthor{video:pit-maven-tutorial} \\
Carga & \citetitle{video:tene-latency} & \citeauthor{video:tene-latency}
  (origem da omissão coordenada) \\
Carga & \citetitle{video:gatling-java-beginners} &
  \citeauthor{video:gatling-java-beginners} \\
Carga & \citetitle{video:gatling-maven-pt} \pt &
  \citeauthor{video:gatling-maven-pt} \\
Caos & \citetitle{video:netflix-principles-chaos} (SREcon17) &
  \citeauthor{video:netflix-principles-chaos}, Netflix \\
Caos & \citetitle{video:rosenthal-chaos-intuition} (GOTO 2016) &
  \citeauthor{video:rosenthal-chaos-intuition} \\
Caos & \citetitle{video:baumgartner-chaos-spring} &
  \citeauthor{video:baumgartner-chaos-spring} \\
Caos & \citetitle{video:fletcher-chaos-monkey-spring} (Spring I/O Bridge) &
  John Fletcher e Manuel Wessner \\
\bottomrule
\end{longtable}
\normalsize

\subsection*{Ferramentas e documentação}
Nesta aula: \citetitle{doc:pitest}, \citetitle{doc:gatling},
\citetitle{doc:jmeter}, \citetitle{doc:k6},
\citetitle{doc:chaos-monkey-spring}, \citetitle{doc:toxiproxy},
\citetitle{doc:netflix-chaosmonkey} e \citetitle{doc:resilience4j}.
Para contratos: \citetitle{doc:pact} (mantido) e
\citetitle{doc:spring-cloud-contract}, que a Spring deixou de manter em
julho de 2026 (o projeto passou para o Stubborn.sh); para integração,
\citetitle{doc:testcontainers}. A lista completa está no apêndice de
ferramentas.
